\documentclass[11pt]{article}

\usepackage[margin=1in]{geometry}
\usepackage{amsmath,amssymb,amsthm,mathtools}
\usepackage{enumitem}
\usepackage{booktabs}
\usepackage{tabularx}
\usepackage{microtype}
\usepackage[hidelinks]{hyperref}
\usepackage{xurl}
\newcolumntype{Y}{>{\raggedright\arraybackslash}X}
\newcolumntype{Z}{>{\raggedright\arraybackslash}p{0.26\linewidth}}
\emergencystretch=4em

\newtheorem{definition}{Definition}
\newtheorem{lemma}{Lemma}
\newtheorem{remark}{Remark}

\title{Public Request Response Packets for Source-Only Releases:\\[0.5ex]
\large Exact Bounded Handoff Slices after Notice-Keyed Follow-Up Choice}
\author{}
\date{Unified archive note v0.13 (2026-06-04; archive rev0837)}

\begin{document}
\maketitle

\begin{abstract}
The source-only branch can now say what a terse paper promises on request, when that promise is complete, whether the paper-level token survives a successor cut, what exact short notice sentence readers saw, what canonical sentence family sits behind that sentence, why that family was selected, and which already-maintained owner is the smallest sufficient answer for one concrete follow-up class.
One practical ambiguity still remained.
Once that owner has been chosen, what exact bounded handoff slice should actually travel?
This note gives that packet cut one home.
It defines compact \emph{public request response packets}: notice-and-follow-up keyed entry-level handoff slices cut from the already-owned source-only objects selected by the response menu.
It also exports one compact first-reopen-owner card so later terse papers can tell which narrower or wider owner regains authority first once one packet shortcut stops being honest.
\end{abstract}

\paragraph{Non-redundancy note.}
Synthesis~56 owns the exact paper-facing source-only contract behind the public sentence, Synthesis~58 owns the family-scoped completion criterion, Synthesis~59 owns the within-release paper-level status token, Synthesis~60 owns the successor carryforward verdict, Synthesis~61 owns the exact outward-facing notice sentence and pointer order, Synthesis~62 owns the canonical sentence family behind that notice, Synthesis~63 owns the token-keyed choice of which such family applies, Synthesis~64 owns the follow-up-class choice of which already-maintained owner is the smallest sufficient answer, Synthesis~66 owns the successor reuse verdict for that exact packet cut, Synthesis~67 owns the packet-local refresh delta, Synthesis~68 owns the exact emitted visible refresh sentence, Synthesis~69 owns the canonical visible wrapper family, Synthesis~70 owns the family-selection owner, Synthesis~71 owns the next successor-facing handoff after that visible layer is fixed, Synthesis~72 owns the exact successor-facing packet cut after that handoff is fixed, Synthesis~73 owns the request-side terminal stop later terse papers should usually cite, Synthesis~74--75 own the paired terminal / cutover discipline, and Synthesis~17 owns the concrete worked instantiation.\cite{series:publicrequestcontracts,series:requestfulfillmentcertificates,series:publicrequeststatusenvelopes,series:publicrequestcarryforward,series:publicrequestnotices,series:publicrequestnoticenormalforms,series:publicrequestnoticeselections,series:publicrequestresponsemenus,series:publicrequestresponsepacketcarryforward,series:publicrequestresponsepacketdeltas,series:publicrequestresponsepacketrefreshnotices,series:publicrequestresponsepacketrefreshnoticenormalforms,series:publicrequestresponsepacketrefreshnoticeselections,series:publicrequestresponsepacketrefreshresponsemenus,series:publicrequestresponsepacketrefreshresponsepackets,series:publicrequestresponsepacketrefreshresponsepacketclosureverdicts,series:pairedterminalcitationdiscipline,series:pairedterminalcutoverrules,series:workedexample}
This note owns only the exact bounded packet cut that travels after one such menu choice has been fixed, not the successor-release rule for whether that packet may later be reused or must reopen.
Later terse papers should cite Synthesis~65 only when the live issue is the exact predecessor source-only response-packet cut itself. They should stop at Synthesis~64 when the live issue is which already-maintained owner is the smallest sufficient answer for one visible notice/follow-up pair, widen to Synthesis~66 when the live issue is whether that exact packet survives the successor release unchanged, refreshed, or reopened, widen to Synthesis~67 when the live issue is which packet-local bindings refreshed, widen to Synthesis~68 when the live issue is the exact emitted visible refresh sentence, widen to Synthesis~69 when the live issue is the canonical visible wrapper family behind that sentence, widen to Synthesis~70 when the live issue is which such family was selected, widen to Synthesis~71 when the live issue is the next successor-facing handoff after that visible layer is fixed, widen to Synthesis~72 when the live issue is the exact successor-facing packet cut after that handoff has already been fixed, widen to Synthesis~73 when the live issue is whether that chosen handoff already closes the request-side source-only branch, drop to Synthesis~17 when one concrete checked-answer/request-tail seam is the live issue instead, and widen to Synthesis~74--75 only when the remaining issue is the abstract paired-terminal discipline or abstract fail-closed cutover.

\paragraph{Scope.}
This is not a new public request contract, not a new completion rule, not a new status token, not a new carryforward verdict, not a new emitted notice, not a new canonical family, and not a new follow-up-class owner choice.
It is a thin packet note answering: after the response menu has chosen the right owner for one concrete follow-up class, what exact bounded slice should actually be handed over?

\paragraph{Exports.}
This note exports (i) public request response packets, (ii) an owner-packet rule, (iii) a follow-up-vocabulary import rule, (iv) a packet-scoping rule, (v) a bounded-companion rule, (vi) a minimality rule saying that later terse papers may cite one maintained packet when they need the exact source-only handoff slice rather than merely the right owner family, and (vii) a compact first-reopen-owner card saying which narrower or wider owner regains authority first once one maintained predecessor packet shortcut goes stale.


\paragraph{One five-stop visible-request route.}
Once a source-only branch has already fixed one emitted notice, one reusable family, one selector, and one follow-up handoff owner, later terse papers should usually traverse one recurring five-stop midbody route: Synthesis~61 for the exact emitted notice, Synthesis~62 for the reusable family, Synthesis~63 for the selector that licenses that family, Synthesis~64 for the smallest-sufficient follow-up handoff, and Synthesis~65 for the exact predecessor-side bounded packet cut that should actually travel. This note is the fifth and last within-release stop before the successor-maintenance tail begins. Stop here when the live issue is the exact bounded carried slice. Hand off left to Synthesis~61--64 when wrapper, family, selector, or handoff choice is still unresolved, and hand off right to Synthesis~66--72 only once the live issue has genuinely become successor reuse, packet-local refresh, visible refresh, successor-facing handoff, or successor-facing packet scope.

\paragraph{A tiny five-stop drill.}
In the worked successor request-scope branch, Synthesis~65 should answer only the question ``what exact contract packet actually travels once \nolinkurl{request_scope_check} has already selected the contract owner?'' It should not also claim to own the visible carryforward sentence, the reusable family, the selector, or the handoff choice itself. Those belong respectively to Synthesis~61, Synthesis~62, Synthesis~63, and Synthesis~64.

\section{Why response packets deserve their own note}
A response menu solves the \emph{which owner?} question.
It does not by itself solve the \emph{which exact bounded cut of that owner should travel?} question.
That distinction matters because the selected owner is often a ledger rather than one scalar field.
A completion-rule check may need exactly the family entries named by the visible notice, plus the public anchor paths and shared validator companions those entries say must travel together.
A request-scope check may need exactly the paper-facing contract entry plus its shared validation companions, not the whole source-only tail.
Without a maintained response-packet note, later papers either overshare a whole imported ledger or improvise a smaller cut in prose.

\begin{definition}[Public request response packet]
Fix one emitted public-request notice $N$ and one follow-up class $c$ covered by the companion response menu $M(N)$.
A \emph{public request response packet}
\[
\Pi(N,c)=(M(N),O,E,P,V)
\]
contains the imported response-menu choice $M(N,c)$, the imported selected owner $O$, one exact entry excerpt $E$ copied from that owner, one set $P$ of required public anchor paths, and one set $V$ of required shared validation companions.
So a response packet is the exact bounded handoff slice for one notice/follow-up pair, not merely the owner family name.
\end{definition}

\begin{definition}[Owner-packet rule]
A public request response packet satisfies the \emph{owner-packet rule} if it imports exactly one response-menu choice for the same notice/follow-up pair and copies only the entry excerpt selected by that choice.
So the packet may not silently replace the response menu, widen to a different owner, or invent a fresh cross-owner bundle.
\end{definition}

\begin{definition}[Follow-up-vocabulary import rule]
A public request response packet satisfies the \emph{follow-up-vocabulary import rule} if its repeated follow-up-class label is copied only from the companion response menu and denotes only that menu's imported shared request-handling vocabulary rather than a packet identity in its own right.
So packetization may reuse the class name without re-owning the follow-up taxonomy.
\end{definition}

\begin{definition}[Packet-scoping rule]
A public request response packet satisfies the \emph{packet-scoping rule} if it is keyed by exactly one emitted notice lineage and one concrete follow-up-class choice from the companion response menu.
So every bounded handoff slice stays attached to one notice/follow-up pair rather than drifting into a notice-independent bundle.
\end{definition}

\begin{definition}[Bounded-companion rule]
A public request response packet satisfies the \emph{bounded-companion rule} if any added paths beyond the selected owner path are limited to:
\begin{enumerate}[leftmargin=*]
\item the imported basis path already named by the response menu,
\item the public anchor paths explicitly required by the selected completion entries, and
\item the shared validation companions already named by the selected contract or completion owner.
\end{enumerate}
So the packet may add only companions that the imported owner already says must travel together.
\end{definition}

\begin{definition}[Minimality rule]
A public request response packet satisfies the \emph{minimality rule} if no packet path comes from an unrelated source-only family once the notice/follow-up choice has been fixed.
So packetization may shrink one owner to its relevant entry slice, but it may not widen to unrelated visible notices, family selectors, or completion families.
\end{definition}

\begin{lemma}[Response packets keep handoff slices explicit without changing source-only ownership]
Assume every public request response packet satisfies the owner-packet rule, the follow-up-vocabulary import rule, the packet-scoping rule, the bounded-companion rule, and the minimality rule.
Then a later terse paper may cite one maintained packet as the exact bounded handoff slice for one concrete follow-up about one visible source-only notice without changing the ownership of request scope, completion, status, carryforward, visible wording, canonical family, or follow-up-class selection.
\end{lemma}

\begin{proof}
The owner-packet rule prevents packetization from changing which already-owned source-only object was selected.
The follow-up-vocabulary import rule makes clear that repeated class labels are only imported menu vocabulary rather than packet identities in their own right.
The packet-scoping rule keeps every bounded handoff slice attached to one concrete notice/follow-up pair.
The bounded-companion rule restricts added paths to companions that the imported owner already requires for auditability.
The minimality rule prevents unrelated families from being smuggled into the handoff once the notice/follow-up pair is fixed.
So the packet changes only packaging granularity: it freezes the exact bounded slice that should travel, while semantics and ownership remain with the imported source-only notes.\qedhere
\end{proof}

\begin{table}[t]
\centering
\small
\begin{tabularx}{\linewidth}{@{}Y Y Y@{}}
\toprule
Follow-up class & Selected owner from Synthesis~64 & Bounded packet cut \\
\midrule
Basis-token check & status/carryforward envelope & envelope entry plus the narrower completion support needed to justify the token \\
Visible-sentence check & emitted notice & notice entry plus the imported basis path it wraps \\
Canonical-family check & notice normal form & family entry plus the emitted notice and imported basis path \\
Family-choice check & notice selection & selector entry plus the chosen normal form and imported basis path \\
Request-scope check & public request contract & contract entry plus the shared validation companions it already names \\
Completion-rule check & request-fulfillment certificates & exactly the covered family entries plus their required public anchor paths and shared validator companions \\
\bottomrule
\end{tabularx}
\caption{A response packet freezes the exact bounded handoff slice after the response menu has already chosen the right owner.}
\label{tab:publicrequestresponsepacket}
\end{table}

\paragraph{Packet stop discipline.}
For the question ``after the visible source-only notice and follow-up-class choice have been fixed, what exact bounded slice should travel?'', stop here. This note owns only that exact bounded packet cut. Under the paired cutover rule, later terse papers should reach this note only after the live question has genuinely widened beyond the answer-side terminal stop into the adjacent source-only branch, and they should continue past this note only if the remaining issue is packet reuse, packet-local delta structure, or the visible refresh wrapper rather than the packet cut itself.\cite{series:pairedterminalcutoverrules} So Synthesis~66--70 remain topic-specific continuations, not part of the default citation path.

\paragraph{A forced public response-packet matrix.}

\begin{table}[t]
\centering
\footnotesize
\begin{tabularx}{\linewidth}{@{}Y Y Y@{}}
\toprule
Worked packet cut & Still the same packet while & Shortcut goes stale when \\
\midrule
\nolinkurl{worked_example_receipt_interlock_successor_public_request_response_menu__request_scope_check} & Same \nolinkurl{public_request_notice_key = worked_example_receipt_interlock_public_request_carryforward_notice}, same \nolinkurl{request_class = request_scope_check}, same \nolinkurl{selected_object_key = worked_example_receipt_interlock_public_request_contract} at \nolinkurl{example_public_request_contracts.json}, same five \nolinkurl{packet_paths}, same empty \nolinkurl{required_public_paths}, and the same four shared \nolinkurl{required_validation_paths}. & Any notice-lineage drift, follow-up-class drift, contract-excerpt drift in the public-anchor/validation floor, or any added or dropped packet path means the same request-scope packet is no longer forced and the packet owner must be reopened. \\
\nolinkurl{worked_example_receipt_interlock_successor_public_request_response_menu__completion_rule_check} & Same \nolinkurl{public_request_notice_key = worked_example_receipt_interlock_public_request_carryforward_notice}, same \nolinkurl{request_class = completion_rule_check}, same \nolinkurl{selected_object_key = worked-example-request-fulfillment-certificates-v1} at \nolinkurl{example_request_fulfillment_certificates.json}, same nine \nolinkurl{packet_paths}, same four \nolinkurl{required_public_paths}, and the same two family-scoped certificate excerpts for \nolinkurl{public_status_sentence_family} and \nolinkurl{compact_diff_summary_family}. & Any drift in the covered family set, the public-anchor floor, the shared validator-companion floor, or the packet cut itself means the same completion-rule packet is no longer forced; the note must reopen Synthesis~65 rather than silently widening into selector/family ledgers or shrinking below the certified family slice. \\
\bottomrule
\end{tabularx}
\caption{A forced public response-packet matrix. Once the visible source-only notice and follow-up class are fixed, the same bounded packet cut remains mechanically forced only while the notice lineage, selected owner excerpt, and companion-path floor remain unchanged.}
\label{tab:forced-public-response-packet-matrix}
\end{table}

This matrix makes the packet stop rule explicit: the same menu entry is not yet enough. A later terse paper gets the same bounded handoff slice only while the exact notice lineage, follow-up label, selected owner excerpt, and carried companion floor still agree. Once any of those drift, the packet shortcut has gone stale and the packet owner must be reopened rather than inferred from the older menu entry alone.\cite{series:workedexample}

\paragraph{One tiny response-packet forcing drill.}
Suppose a later terse note still answers the successor carryforward notice's \nolinkurl{request_scope_check} follow-up and still carries exactly the contract excerpt together with the same five request-scope packet paths. Then the same packet \nolinkurl{worked_example_receipt_interlock_successor_public_request_response_menu__request_scope_check} is still forced. The moment the note adds \nolinkurl{example_compare_report.json} or \nolinkurl{example_successor_clause_pack.json}, it has already widened to the completion-rule packet instead, and the smaller contract packet is stale. If the notice lineage changes or the selected owner stops being the contract owner, the note must reopen Synthesis~65 rather than pretending that the old request-scope packet still travels.\cite{series:workedexample}

\paragraph{A minimal public response-packet card.}

\begin{table}[t]
\centering
\footnotesize
\begin{tabularx}{\linewidth}{@{}Z Y@{}}
\toprule
Card field & Worked floor \\
\midrule
Smallest nontrivial packet keys & \nolinkurl{worked_example_receipt_interlock_within_release_public_request_response_menu__request_scope_check} and \nolinkurl{worked_example_receipt_interlock_successor_public_request_response_menu__request_scope_check} \\
Selected owner / scope floor & Both smallest packets select the single contract owner \nolinkurl{worked_example_receipt_interlock_public_request_contract} at \nolinkurl{example_public_request_contracts.json} and keep \nolinkurl{packet_scope_mode = notice_and_followup_scoped_packet} \\
Smallest packet cut floor & Both request-scope packets freeze exactly five packet paths: \nolinkurl{../tools/validate_example.py}, \nolinkurl{example_artifact_inventory.json}, \nolinkurl{example_public_request_contracts.json}, \nolinkurl{example_validation_report.json}, and \nolinkurl{support_manifest.json} \\
Smallest packet anchor floor & Both request-scope packets leave \nolinkurl{required_public_paths} empty and place only the four shared maintenance companions in \nolinkurl{required_validation_paths} \\
Largest bounded packet keys & \nolinkurl{worked_example_receipt_interlock_within_release_public_request_response_menu__completion_rule_check} and \nolinkurl{worked_example_receipt_interlock_successor_public_request_response_menu__completion_rule_check} \\
Largest packet cut floor & Both completion-rule packets freeze exactly nine packet paths: the certificate ledger \nolinkurl{example_request_fulfillment_certificates.json}, four public anchors \nolinkurl{example_compare_report.json}, \nolinkurl{example_successor_claim_support_map.json}, \nolinkurl{example_successor_clause_pack.json}, \nolinkurl{example_successor_lineage_notice.json}, and the same four validation companions \\
Largest packet excerpt floor & Both completion-rule packets carry exactly the two family-scoped certificate excerpts for \nolinkurl{public_status_sentence_family} and \nolinkurl{compact_diff_summary_family}; they do not widen to unrelated notice-family or selector owners once the follow-up class is fixed \\
Escalation pointer & Stop at Synthesis~64 when the live issue is still which owner answers the follow-up; widen to Synthesis~66--72 only when successor reuse / delta / visible refresh / successor-facing packet ownership is already live; widen to Synthesis~73 for the maintained request-side terminal stop, drop to Synthesis~17 for one concrete checked-answer/request-tail seam, and widen to Synthesis~74--75 only for the abstract paired-tail discipline or abstract cutover judgment \\
\bottomrule
\end{tabularx}
\caption{A minimal public response-packet card. This is the smallest public answer later terse papers can quote directly when the live issue is the exact bounded handoff slice that should travel for one notice/follow-up pair rather than merely which maintained owner was selected or whether a successor refresh branch later reopens.}
\label{tab:minimal-public-response-packet-card}
\end{table}

This card is intentionally non-decorative: it pins the actual worked request-scope and completion-rule packet keys, the exact packet-path counts, the explicit public-anchor versus validation-companion split, and the family-local certificate floor without silently re-owning the response menu, the successor carryforward rules, or the request-side terminal closure verdict.\cite{series:workedexample}

\paragraph{One tiny response-packet drill.}
Suppose a later terse note has already fixed the successor menu entry \nolinkurl{worked_example_receipt_interlock_successor_public_request_response_menu__request_scope_check}. The honest packet cut is then the five-path request-scope packet and nothing larger: the contract path plus its four shared maintenance companions. If the live question instead shifts to whether the promised families are complete, the honest packet changes to \nolinkurl{worked_example_receipt_interlock_successor_public_request_response_menu__completion_rule_check}, which widens exactly by the four public anchor files and the two family-scoped certificate excerpts. In neither case may a later paper quietly add the notice-normal-form or selector ledgers once the follow-up class is fixed, because that would stop being a packet cut and would become a wider source-only rewalk.\cite{series:workedexample}

\paragraph{A tiny response-packet sentence card.}
\begin{table}[t]
\centering
\small
\begin{tabularx}{\linewidth}{@{}Z Y@{}}
\toprule
Field & Minimal public answer \\
\midrule
Packet clause floor & \texttt{Response packet }\emph{[packet key]}\texttt{ cuts }\emph{[excerpt key]}\texttt{ from }\emph{[owner path]}\texttt{ for notice }\emph{[notice key]}\texttt{ and follow-up class }\emph{[request class]}\texttt{, carrying exactly }\emph{[packet-path phrase]}\texttt{ with public-anchor floor }\emph{[public-anchor phrase]}\texttt{ and validation-companion floor }\emph{[validation phrase]}. \\
Imported-owner floor & The packet, excerpt, owner path, notice, and follow-up-class slots remain the already-owned packet ledger and menu choice; the clause does not paraphrase them away into a fresh owner summary \\
Cut floor & The packet-path phrase is the exact bounded travelled cut, not a looser gesture toward ``supporting files'' that silently adds or drops paths after the follow-up class has already fixed the owner \\
Companion-floor discipline & The public-anchor and validation-companion slots must preserve the packet's explicit split between outward public anchors and shared maintenance companions; they may not be fused into one vague ``evidence'' bucket \\
Owner-visibility floor & The sentence keeps eight narrow anchors visible at once: packet key, excerpt key, owner path, notice key, request class, packet-path phrase, public-anchor phrase, and validation-companion phrase \\
Escalation pointer & Stop at Synthesis~64 when the live issue is still which owner the notice/follow-up pair selects, stop at Synthesis~66 when the live issue has moved from the predecessor packet cut to successor reuse discipline, and then widen to Synthesis~73 for the maintained request-side terminal stop, drop to Synthesis~17 for one concrete checked-answer/request-tail seam, and widen to Synthesis~74--75 only for the abstract paired-tail discipline or abstract cutover judgment rather than the predecessor packet itself \\
\bottomrule
\end{tabularx}
\caption{A tiny response-packet sentence card. This is the smallest sentence-level rendering later terse papers can quote directly when the live issue is the exact predecessor-side bounded packet cut itself rather than merely the owner choice or the later successor-maintenance tail.}
\label{tab:tiny-response-packet-sentence-card}
\end{table}

This sentence card is intentionally non-decorative: it turns the already-owned predecessor packet cut into one exact clause for later terse papers, but it does not re-own the response-menu choice, the successor reuse verdict, the packet-local delta ledger, or the downstream request-side terminal / paired cutover layers that neighboring papers already own.\cite{series:workedexample}

\paragraph{One tiny response-packet sentence fill.}
For the worked successor request-scope branch, the sentence exported by this note is simply: \emph{Response packet \path{worked_example_receipt_interlock_successor_public_request_response_menu__request_scope_check} cuts \path{worked_example_receipt_interlock_public_request_contract} from \path{example_public_request_contracts.json} for notice \path{worked_example_receipt_interlock_public_request_carryforward_notice} and follow-up class \path{request_scope_check}, carrying exactly \path{../tools/validate_example.py}, \path{example_artifact_inventory.json}, \path{example_public_request_contracts.json}, \path{example_validation_report.json}, and \path{support_manifest.json} with public-anchor floor empty and validation-companion floor equal to the four shared maintenance companions.} If any one of those eight owned anchors changes, later terse papers should reopen the response-packet owner instead of silently keeping the old clause. If the live issue is still which owner was selected or has already shifted to the successor reuse verdict, they should fall back to Synthesis~64 or widen to Synthesis~66 rather than stretching this predecessor packet sentence beyond packet-cut scope.\cite{series:workedexample}


\paragraph{A compact first-reopen-owner card.}
\begin{table}[t]
\centering
\small
\begin{tabularx}{\linewidth}{@{}Z Y@{}}
\toprule
Field & Minimal routing answer \\
\midrule
Local packet-cut drift & Reopen \textbf{Synthesis~65} when the packet key, excerpt key, owner path, packet-path cut, required public-anchor set, or required validation-companion set changes while the surrounding request-side visible branch is otherwise still recognizable \\
Response-menu handoff drift & Reopen \textbf{Synthesis~64} when the notice/follow-up pair no longer selects the same maintained owner or the same predecessor packet can no longer honestly claim to be copied from the same response-menu choice \\
Within-release basis drift & Reopen \textbf{Synthesis~59} when the imported within-release status envelope or its paper-level token no longer matches the basis this predecessor packet branch imports \\
Successor carryforward basis drift & Reopen \textbf{Synthesis~60} when the imported successor carryforward envelope, release-pair status token, or carryforward verdict no longer matches the basis this predecessor packet branch imports \\
Emitted visible-wrapper drift & Reopen \textbf{Synthesis~61} when the packet still points at the same branch but the exact outward-facing sentence or pointer order ceases to agree with the emitted notice key that the response menu and packet import \\
Canonical-family drift & Reopen \textbf{Synthesis~62} when the packet still travels for the same visible branch but the reusable canonical family, compatible-token set, slot bindings, or render map behind that notice has changed \\
Notice-selector drift & Reopen \textbf{Synthesis~63} when the imported token-to-family choice or selected emitted-notice lineage no longer matches the packet branch this note imports \\
Successor packet-maintenance remainder & Widen to \textbf{Synthesis~66--72} when the predecessor packet cut is fixed and the live issue has moved to successor reuse, packet-local delta, visible refresh, canonical refresh family, refresh selection, successor response-menu choice, or successor packet scope \\
Terminal / paired remainder & Widen to \textbf{Synthesis~73} for the maintained request-side terminal stop, drop to \textbf{Synthesis~17} for one concrete checked-answer/request-tail seam, and widen to \textbf{Synthesis~74--75} only for the abstract paired-terminal discipline or abstract cutover judgment when the downstream request-side branch is already fixed \\
\bottomrule
\end{tabularx}
\caption{A compact first-reopen-owner card. This is the smallest local map from one stale public-request-response-packet shortcut to the first narrower or wider owner that regains authority.}
\label{tab:first-reopen-owner-card-public-request-response-packets}
\end{table}

\paragraph{One tiny first-reopen-owner drill.}
If a later terse paper still asks \emph{``what exact bounded source-only slice should actually travel for this already-chosen visible follow-up?''} while the surrounding request-side visible branch remains recognizable but the packet key, excerpt key, owner path, packet-path cut, public-anchor set, or validation-companion set changes, Table~\ref{tab:first-reopen-owner-card-public-request-response-packets} sends the paper back to Synthesis~65 itself. If the same predecessor packet can no longer honestly claim to be copied from the same notice/follow-up response-menu choice, the same table reopens Synthesis~64. If the packet branch is fixed but the imported within-release status basis drifts, it reopens Synthesis~59; if instead the imported successor carryforward basis drifts, it reopens Synthesis~60. If the exact outward-facing sentence or pointer order changes, it reopens Synthesis~61; if the reusable canonical family, compatible-token set, slot bindings, or render map behind that notice changes, it reopens Synthesis~62; if the imported token-to-family choice or selected emitted-notice lineage changes, it reopens Synthesis~63. If the predecessor packet cut is already fixed and the real question has moved to successor reuse, packet-local delta, visible refresh, canonical refresh family, refresh selection, successor response-menu choice, or successor packet scope, the same table widens to Synthesis~66--72. If that downstream request-side branch is also fixed and only the request-side terminal stop, one concrete checked-answer/request-tail seam, or the abstract paired cutover remains live, it widens to Synthesis~73 for the maintained request-side terminal stop, drops to Synthesis~17 for one concrete seam, and widens to Synthesis~74--75 only for the abstract paired-terminal discipline or abstract cutover judgment instead of pretending that the predecessor packet cut still owns the remaining judgment.\cite{series:publicrequeststatusenvelopes,series:publicrequestcarryforward,series:publicrequestnotices,series:publicrequestnoticenormalforms,series:publicrequestnoticeselections,series:publicrequestresponsemenus,series:publicrequestresponsepacketcarryforward,series:publicrequestresponsepacketdeltas,series:publicrequestresponsepacketrefreshnotices,series:publicrequestresponsepacketrefreshnoticenormalforms,series:publicrequestresponsepacketrefreshnoticeselections,series:publicrequestresponsepacketrefreshresponsemenus,series:publicrequestresponsepacketrefreshresponsepackets,series:publicrequestresponsepacketrefreshresponsepacketclosureverdicts,series:pairedterminalcitationdiscipline,series:pairedterminalcutoverrules}

\section{Worked example pointer}
The maintained worked bundle now ships \nolinkurl{example_public_request_response_packets.json}.\cite{series:workedexample}
Its worked packet ledger is now concrete enough that later terse papers can cite exact packet keys and exact travelled file cuts instead of describing them generically:
\begin{enumerate}[leftmargin=*]
\item the ledger contains exactly twelve packet entries, obtained by crossing the two worked menu keys \nolinkurl{worked_example_receipt_interlock_within_release_public_request_response_menu} and \nolinkurl{worked_example_receipt_interlock_successor_public_request_response_menu} with the six shared follow-up classes \nolinkurl{basis_token_check}, \nolinkurl{visible_sentence_check}, \nolinkurl{canonical_family_check}, \nolinkurl{family_choice_check}, \nolinkurl{request_scope_check}, and \nolinkurl{completion_rule_check};
\item the two \nolinkurl{request_scope_check} packets are the smallest nontrivial contract packets: both select \nolinkurl{worked_example_receipt_interlock_public_request_contract} at \nolinkurl{example_public_request_contracts.json}, both freeze the same five packet paths \nolinkurl{../tools/validate_example.py}, \nolinkurl{example_artifact_inventory.json}, \nolinkurl{example_public_request_contracts.json}, \nolinkurl{example_validation_report.json}, and \nolinkurl{support_manifest.json}, and both leave \nolinkurl{required_public_paths} empty while placing only the four maintenance companions in \nolinkurl{required_validation_paths}; in particular, the first adjacent source-only widening step beneath the answer-side compact-diff review drill is now concrete rather than implied: \nolinkurl{example_public_request_notice.json} $\to$ \nolinkurl{worked_example_receipt_interlock_successor_public_request_response_menu} $\to$ \nolinkurl{worked_example_receipt_interlock_successor_public_request_response_menu__request_scope_check};
\item the two \nolinkurl{completion_rule_check} packets are the largest bounded certificate packets: both select \nolinkurl{worked-example-request-fulfillment-certificates-v1} at \nolinkurl{example_request_fulfillment_certificates.json}, both add exactly the four public companions \nolinkurl{example_compare_report.json}, \nolinkurl{example_successor_claim_support_map.json}, \nolinkurl{example_successor_clause_pack.json}, and \nolinkurl{example_successor_lineage_notice.json}, and both reuse the same four maintenance validation companions; and
\item the maintained question-route catalog mirrors these exact packet keys in \nolinkurl{source_only_response_packet_map}, carries the packet support-surface cuts separately in \nolinkurl{source_only_response_packet_inventory_cut_map}, \nolinkurl{source_only_response_packet_manifest_cut_map}, and \nolinkurl{source_only_response_packet_support_surface_map}, and states in \nolinkurl{source_only_response_packet_scope_mode_map} that all twelve entries are \nolinkurl{notice_and_followup_scoped_packet}.
\end{enumerate}
So a later terse note can recover both the semantic packet identity and the exact travelled file cut from one maintained route object without pretending that the route itself owns packet semantics. Under the paired terminal citation discipline, later terse papers should usually cite Synthesis~73 for the maintained request-side terminal stop, drop to Synthesis~17 when they want one concrete checked-answer/request-tail seam, and widen to Synthesis~74--75 only for the abstract paired-tail discipline or abstract cutover judgment rather than reopening this packet note; they should return here only when the exact packet cut is itself the live issue.\cite{series:pairedterminalcitationdiscipline,series:pairedterminalcutoverrules} If a later paper instead needs the successor-release reuse discipline for those exact packets, that narrower carryforward owner belongs in Synthesis~66 rather than here.

\begin{thebibliography}{99}\small

\bibitem{series:publicrequestcontracts}
Anonymous.
\newblock Public Request Contracts for Source-Only Releases: Exact Paper-Level Handoffs from Boilerplate Policy to Maintained Packet Cuts.
\newblock Series synthesis note (Synthesis~56), 2026. (This archive.)

\bibitem{series:requestfulfillmentcertificates}
Anonymous.
\newblock Request Fulfillment Certificates for Source-Only Releases: Exact Completion Criteria for Public Request Contracts and Typed Omitted-Support Cuts.
\newblock Series synthesis note (Synthesis~58), 2026. (This archive.)

\bibitem{series:publicrequeststatusenvelopes}
Anonymous.
\newblock Public Request Status Envelopes for Source-Only Releases: Paper-Level Tokens Compressing Family Completion without Re-owning the Promise.
\newblock Series synthesis note (Synthesis~59), 2026. (This archive.)

\bibitem{series:publicrequestcarryforward}
Anonymous.
\newblock Public Request Carryforward for Source-Only Successor Releases: When Paper-Level Request Status Persists, Advances, or Reopens across Release Evolution.
\newblock Series synthesis note (Synthesis~60), 2026. (This archive.)

\bibitem{series:publicrequestnotices}
Anonymous.
\newblock Public Request Notices for Source-Only Releases: Outward-Facing Sentences for Within-Release Status and Cross-Release Carryforward.
\newblock Series synthesis note (Synthesis~61), 2026. (This archive.)

\bibitem{series:publicrequestnoticenormalforms}
Anonymous.
\newblock Public Request Notice Normal Forms for Source-Only Releases: Canonical Sentence Families for Within-Release Status and Cross-Release Carryforward Notices.
\newblock Series synthesis note (Synthesis~62), 2026. (This archive.)

\bibitem{series:publicrequestnoticeselections}
Anonymous.
\newblock Public Request Notice Selections for Source-Only Releases: Choosing Canonical Notice Families from Within-Release Status and Cross-Release Carryforward Tokens.
\newblock Series synthesis note (Synthesis~63), 2026. (This archive.)

\bibitem{series:publicrequestresponsemenus}
Anonymous.
\newblock Public Request Response Menus for Source-Only Releases: Smallest-Sufficient Follow-Up Owners for Visible Notices.
\newblock Series synthesis note (Synthesis~64), 2026. (This archive.)

\bibitem{series:publicrequestresponsepacketcarryforward}
Anonymous.
\newblock Public Request Response Packet Carryforward for Source-Only Successor Releases: Reuse Rules for Exact Notice-Keyed Handoff Slices under Release Evolution.
\newblock Series synthesis note (Synthesis~66), 2026. (This archive.)

\bibitem{series:publicrequestresponsepacketdeltas}
Anonymous.
\newblock Public Request Response Packet Delta Ledgers for Source-Only Successor Releases: Exact Refresh Instructions for Reused Notice-Keyed Handoff Slices.
\newblock Series synthesis note (Synthesis~67), 2026. (This archive.)

\bibitem{series:publicrequestresponsepacketrefreshnotices}
Anonymous.
\newblock Public Request Response Packet Refresh Notices for Source-Only Successor Releases: Exact Visible Reissue Sentences for Reused Notice-Keyed Handoff Slices.
\newblock Series synthesis note (Synthesis~68), 2026. (This archive.)

\bibitem{series:publicrequestresponsepacketrefreshnoticenormalforms}
Anonymous.
\newblock Public Request Response Packet Refresh Notice Normal Forms for Source-Only Successor Releases: Canonical Sentence Families for Visible Reissue Wrappers of Reused Notice-Keyed Handoff Slices.
\newblock Series synthesis note (Synthesis~69), 2026. (This archive.)

\bibitem{series:publicrequestresponsepacketrefreshnoticeselections}
Anonymous.
\newblock Public Request Response Packet Refresh Notice Selections for Source-Only Successor Releases: Choosing Canonical Visible Reissue Families from Packet-Local Delta Tokens.
\newblock Series synthesis note (Synthesis~70), 2026. (This archive.)

\bibitem{series:publicrequestresponsepacketrefreshresponsemenus}
Anonymous.
\newblock Public Request Response Packet Refresh Response Menus for Source-Only Successor Releases: Smallest-Sufficient Follow-Up Handoffs for Packet Membership, Reuse Verdicts, Delta Ledgers, Visible Reissue Sentences, Canonical Families, and Family Choice.
\newblock Series synthesis note (Synthesis~71), 2026. (This archive.)

\bibitem{series:publicrequestresponsepacketrefreshresponsepackets}
Anonymous.
\newblock Public Request Response Packet Refresh Response Packets for Source-Only Successor Releases: Exact Bounded Handoff Slices after Concrete Follow-Up Choice about Visible Reissue Sentences.
\newblock Series synthesis note (Synthesis~72), 2026. (This archive.)

\bibitem{series:publicrequestresponsepacketrefreshresponsepacketclosureverdicts}
Anonymous.
\newblock Public Request Response Packet Refresh Response Packet Closure Verdicts for Source-Only Successor Releases: Certifying Stable Finite Basis after Exact Successor-Facing Handoff Slice Choice.
\newblock Series synthesis note (Synthesis~73), 2026. (This archive.)

\bibitem{series:pairedterminalcitationdiscipline}
Anonymous.
\newblock Paired Terminal Citation Discipline for Stable Review Spines: One Answer-Side Stop and One Request-Side Capstone for Terse Papers.
\newblock Series synthesis note (Synthesis~74), 2026. (This archive.)

\bibitem{series:pairedterminalcutoverrules}
Anonymous.
\newblock Paired Terminal Cutover Rules for Stable Review Spines: When to Stop at the Checked Answer, When to Widen to Source-Only, and When to Fail Closed to the Reopening Owner.
\newblock Series synthesis note (Synthesis~75), 2026. (This archive.)

\bibitem{series:workedexample}
Anonymous.
\newblock Worked Example: Receipt Interlock for an Anonymous-DHT Reader-Privacy Claim.
\newblock Series synthesis note (Synthesis~17), 2026. (This archive.)

\end{thebibliography}
\end{document}
